\section{Introducción: tendencias de la GenAI desde la perspectiva empresarial}

Burak Gokturk es vicepresidente de Ingeniería de Google. Esta clase no profundiza en detalles algorítmicos, sino que examina las tendencias de desarrollo de la GenAI desde la perspectiva macro del \textbf{mercado empresarial}, para ofrecer a estudiantes y emprendedores una orientación práctica.

\begin{knowledgebox}{El cambio de perspectiva del doctorado a la industria}
Burak Gokturk abre con su propia experiencia de doctorado: en aquella época se consideraba que las redes neuronales «no eran tan buenas como las SVM», y entrenar los datos de tomografía computarizada de 50 pacientes en 42 segundos ya se percibía como demasiado lento. Hoy en día, entrenar los modelos grandes toma semanas enteras. Una preocupación interesante es que los ciclos de entrenamiento tan prolongados han hecho más lenta la velocidad con la que los investigadores adquieren intuición.
\end{knowledgebox}

\subsection{Resumen de la sección}

La velocidad de desarrollo de la IA supera con creces lo esperado, y tanto la academia como la industria están viviendo un cambio de paradigma.

